\documentclass[12pt]{article}
\usepackage[utf8]{inputenc}
\usepackage{amsmath}
\usepackage{amssymb}
\usepackage{graphicx}
\usepackage{booktabs}
\usepackage{listings}
\usepackage{xcolor}
\usepackage{hyperref}
\usepackage{geometry}
\geometry{margin=1in}

\lstset{
    language=Python,
    basicstyle=\ttfamily\small,
    keywordstyle=\color{blue}\bfseries,
    commentstyle=\color{green!60!black},
    stringstyle=\color{red},
    numbers=left,
    numberstyle=\tiny\color{gray},
    stepnumber=1,
    numbersep=5pt,
    backgroundcolor=\color{gray!10},
    frame=single,
    breaklines=true,
    breakatwhitespace=false,
    tabsize=4,
    showstringspaces=false
}

\title{Benefit-Cost Ratio Calculator Documentation}
\author{}
\date{}

\begin{document}

\maketitle

\section{Introduction}

The \texttt{bcr\_calculator.py} script calculates Benefit-Cost Ratios (BCR) and Net Benefits for transmission projects. It aggregates benefits and costs from all CTCC modules and computes multiple BCR metrics to evaluate project economics from different perspectives.

\section{Method Overview}

\subsection{Purpose}

This script calculates BCR metrics by:
\begin{enumerate}
    \item Loading scenario data from \texttt{batch\_summary.csv}
    \item Aggregating benefits (congestion, curtailment, line loss benefits for reconductoring)
    \item Aggregating costs by category (capital, operational, risk, delay)
    \item Calculating multiple BCR metrics (system, capital, excluding risk, excluding emissions)
    \item Computing net benefits (total and excluding risk)
    \item Displaying formatted summary and updating CSV
\end{enumerate}

\subsection{Calculation Approach}

The BCR calculation follows this structure:

\begin{align}
\text{Total Benefits (Conservative)} &= \text{Congestion (Haircut)} + \text{Curtailment (Haircut)} + \text{Line Loss Benefits} \\
\text{Total Costs} &= \text{Capital} + \text{Operational} + \text{Risk} + \text{Delay} \\
\text{BCR}_{system} &= \frac{\text{Conservative Benefits}}{\text{Total Costs}} \\
\text{BCR}_{capital} &= \frac{\text{Conservative Benefits}}{\text{Capital Costs}} \\
\text{BCR}_{excluding\_risk} &= \frac{\text{Conservative Benefits}}{\text{Total Costs} - \text{Risk Costs}} \\
\text{BCR}_{excluding\_emissions} &= \frac{\text{Conservative Benefits}}{\text{Total Costs} - \text{Emissions Costs}} \\
\text{Net Benefit} &= \text{Conservative Benefits} - \text{Total Costs}
\end{align}

\subsection{Inputs and Outputs}

\textbf{Inputs:}
\begin{itemize}
    \item Scenario ID (to load from \texttt{batch\_summary.csv})
    \item All cost and benefit metrics from CTCC modules
\end{itemize}

\textbf{Outputs:}
\begin{itemize}
    \item BCR metrics (system, capital, excluding risk, excluding emissions) - all use conservative benefits
    \item Net benefits (total and excluding risk)
    \item Benefit and cost breakdowns
    \item Formatted summary display
\end{itemize}

\subsection{Integration with Other Scripts}

This script integrates with:
\begin{itemize}
    \item \texttt{csv\_output\_manager.py} - Reads from and writes to \texttt{batch\_summary.csv} (see \texttt{csv\_output\_manager\_documentation.tex})
    \item All CTCC cost/benefit calculation scripts - Aggregates their outputs
\end{itemize}

\section{Key Functions}

\subsection{\texttt{load\_scenario\_data(scenario\_id, output\_dir)}}

Loads scenario data from \texttt{batch\_summary.csv} for the given scenario ID.

\textbf{Parameters:}
\begin{itemize}
    \item \texttt{scenario\_id} (str): Unique identifier for the scenario
    \item \texttt{output\_dir} (str): Directory containing \texttt{batch\_summary.csv}
\end{itemize}

\textbf{Returns:}
\begin{itemize}
    \item Dictionary with all data for the scenario, or None if not found
\end{itemize}

\textbf{Key Logic:}
\begin{lstlisting}
# Read CSV and find matching scenario_id
with open(batch_path, 'r') as f:
    reader = csv.DictReader(f)
    for row in reader:
        if row.get('scenario_id') == scenario_id:
            # Convert numeric strings to floats
            for key, value in row.items():
                if value and value != '':
                    try:
                        row[key] = float(value)
                    except (ValueError, TypeError):
                        pass
            return row
\end{lstlisting}

\subsection{\texttt{calculate\_benefits(data)}}

Calculates total benefits from scenario data.

\textbf{Parameters:}
\begin{itemize}
    \item \texttt{data} (dict): Dictionary with scenario data
\end{itemize}

\textbf{Returns:}
\begin{itemize}
    \item Dictionary with benefit breakdown and totals (PV, nominal, haircut)
\end{itemize}

\textbf{Key Logic:}
\begin{lstlisting}
# Present value benefits
congestion_benefit_pv = data.get('congestion_benefit_pv', 0) or 0
curtailment_benefit_pv = data.get('curtailment_benefit_pv', 0) or 0
line_loss_pv = data.get('line_loss_cost_pv', 0) or 0

# For reconductoring, line losses are negative (benefit)
line_loss_benefit_pv = 0
if line_loss_pv < 0:
    line_loss_benefit_pv = abs(line_loss_pv)

total_benefits_pv = (congestion_benefit_pv + curtailment_benefit_pv 
                     + line_loss_benefit_pv)

# Haircut benefits (conservative)
total_benefits_haircut_pv = (congestion_benefit_haircut 
                            + curtailment_benefit_haircut 
                            + line_loss_benefit_pv)
\end{lstlisting}

\subsection{\texttt{calculate\_costs(data)}}

Calculates total costs from scenario data, organized by category.

\textbf{Parameters:}
\begin{itemize}
    \item \texttt{data} (dict): Dictionary with scenario data
\end{itemize}

\textbf{Returns:}
\begin{itemize}
    \item Dictionary with cost breakdown by category and totals (PV and nominal)
\end{itemize}

\textbf{Cost Categories:}
\begin{itemize}
    \item \textbf{Capital}: Build, ROW, environmental mitigation
    \item \textbf{Operational}: O\&M, insurance, line losses (greenfield), emissions
    \item \textbf{Risk}: Wildfire, outage
    \item \textbf{Delay}: Construction delay, congestion/curtailment delay, residual congestion
\end{itemize}

\textbf{Key Logic:}
\begin{lstlisting}
# Capital costs
capital_costs_pv = build_cost_pv + row_cost_pv + env_mitigation_pv

# Operational costs (line losses only if positive - greenfield)
line_loss_cost_pv = max(0, line_loss_pv)
operational_costs_pv = oandm_pv + insurance_pv + line_loss_cost_pv + emissions_pv

# Risk costs
risk_costs_pv = wildfire_pv + outage_pv

# Delay costs
delay_costs_pv = (delay_cost_pv + congestion_delay_pv 
                 + curtailment_delay_pv + residual_congestion_pv)

# Totals
total_costs_pv = (capital_costs_pv + operational_costs_pv 
                 + risk_costs_pv + delay_costs_pv)
\end{lstlisting}

\subsection{\texttt{calculate\_bcr\_metrics(benefits, costs)}}

Calculates benefit-cost ratios and net benefits.

\textbf{Parameters:}
\begin{itemize}
    \item \texttt{benefits} (dict): Dictionary with benefit breakdown
    \item \texttt{costs} (dict): Dictionary with cost breakdown
\end{itemize}

\textbf{Returns:}
\begin{itemize}
    \item Dictionary with BCR metrics and net benefits
\end{itemize}

\textbf{BCR Metrics:}
\begin{itemize}
    \item \textbf{BCR\_system}: Conservative benefits / Total costs (primary metric, uses saturation-adjusted benefits)
    \item \textbf{BCR\_capital}: Conservative benefits / Capital costs (investor perspective, uses saturation-adjusted benefits)
    \item \textbf{BCR\_excluding\_risk}: Conservative benefits / (Total costs - Risk costs) (core economics, uses saturation-adjusted benefits)
    \item \textbf{BCR\_excluding\_emissions}: Conservative benefits / (Total costs - Emissions costs) (uses saturation-adjusted benefits)
\end{itemize}

\textbf{Key Logic:}
\begin{lstlisting}
# Calculate costs excluding risk
total_costs_excluding_risk_pv = total_costs_pv - risk_costs_pv

# Use conservative (haircut) benefits for all BCR calculations
total_benefits_pv = total_benefits_haircut_pv

# Calculate BCRs (with division by zero protection)
bcr_system = total_benefits_pv / total_costs_pv if total_costs_pv > 0 else 0
bcr_capital = total_benefits_pv / capital_costs_pv if capital_costs_pv > 0 else 0
bcr_excluding_risk = (total_benefits_pv / total_costs_excluding_risk_pv 
                      if total_costs_excluding_risk_pv > 0 else 0)
bcr_excluding_emissions = (total_benefits_pv / total_costs_excluding_emissions_pv 
                           if total_costs_excluding_emissions_pv > 0 else 0)

# Net benefits (using conservative benefits)
net_benefit_pv = total_benefits_pv - total_costs_pv
net_benefit_excluding_risk_pv = total_benefits_pv - total_costs_excluding_risk_pv
\end{lstlisting}

\subsection{\texttt{calculate\_and\_display\_bcr(scenario\_id, output\_dir)}}

Main function to calculate and display BCR analysis.

\textbf{Parameters:}
\begin{itemize}
    \item \texttt{scenario\_id} (str): Unique identifier for the scenario
    \item \texttt{output\_dir} (str): Directory containing \texttt{batch\_summary.csv}
\end{itemize}

\textbf{Returns:}
\begin{itemize}
    \item Dictionary with all BCR results, or None if calculation fails
\end{itemize}

\subsection{\texttt{print\_bcr\_summary(benefits, costs, bcr\_metrics, data)}}

Prints formatted BCR summary to terminal.

\section{Example}

The following example demonstrates a typical usage:

\begin{lstlisting}
# Example: 500 MW Overhead AC line
# Scenario ID: "20250101_120000"
# 
# Inputs (from batch_summary.csv):
#   - Congestion benefit PV: $500M
#   - Curtailment benefit PV: $200M
#   - Line loss cost PV: $10M (greenfield - cost)
#   - Build cost PV: $100M
#   - ROW cost PV: $20M
#   - Environmental PV: $15M
#   - O&M PV: $50M
#   - Insurance PV: $5M
#   - Emissions PV: $8M
#   - Wildfire PV: $1B
#   - Outage PV: $500M
#   - Delay costs PV: $2M
#   - Residual congestion PV: $50M
#
# Calculation:
#   Total Benefits = $500M + $200M + $0 = $700M
#   Capital Costs = $100M + $20M + $15M = $135M
#   Operational Costs = $50M + $5M + $10M + $8M = $73M
#   Risk Costs = $1B + $500M = $1.5B
#   Delay Costs = $2M + $50M = $52M
#   Total Costs = $135M + $73M + $1.5B + $52M = $1.76B
#   
#   BCR_system = $700M / $1.76B = 0.40
#   BCR_capital = $700M / $135M = 5.19
#   BCR_excluding_risk = $700M / ($1.76B - $1.5B) = 2.69
#   Net Benefit = $700M - $1.76B = -$1.06B
#
# Output:
#   - BCR_system: 0.40
#   - BCR_capital: 5.19
#   - BCR_excluding_risk: 2.69
#   - Net benefit: -$1.06B
\end{lstlisting}

\textbf{Integration Example:}

\begin{lstlisting}
# When called from ctcc.py:
from scripts.bcr_calculator import calculate_and_display_bcr

# After all other scripts complete
scenario_id = os.environ.get('CTCC_SCENARIO_ID')
results = calculate_and_display_bcr(scenario_id)

# Update batch_summary.csv with BCR metrics
csv_manager = CTCCOutputManager()
csv_manager.add_bcr_metrics(results)
csv_manager.write_batch_summary()
\end{lstlisting}

\section{Key Implementation Details}

\subsection{Line Loss Treatment}

The script handles line losses differently for greenfield vs. reconductoring:
\begin{itemize}
    \item \textbf{Greenfield}: Line losses are positive (costs) - counted in operational costs
    \item \textbf{Reconductoring}: Line losses are negative (benefits) - counted in benefits
\end{itemize}

\subsection{Multiple BCR Perspectives}

The script calculates multiple BCR metrics to provide different perspectives. All BCR metrics use conservative (saturation-adjusted) benefits to account for diminishing marginal returns:
\begin{itemize}
    \item \textbf{System BCR}: Overall project economics including all costs (uses conservative benefits)
    \item \textbf{Capital BCR}: Benefits relative to capital investment (investor view, uses conservative benefits)
    \item \textbf{Excluding Risk BCR}: Core economics without wildfire/outage risk costs (uses conservative benefits)
    \item \textbf{Excluding Emissions BCR}: Economics without emissions costs (uses conservative benefits)
\end{itemize}

\subsection{Net Benefits}

The script calculates net benefits in addition to BCRs:
\begin{itemize}
    \item \textbf{Net Benefit}: Total benefits - Total costs
    \item \textbf{Net Benefit Excluding Risk}: Total benefits - (Total costs - Risk costs)
\end{itemize}

\subsection{Division by Zero Protection}

All BCR calculations include checks to prevent division by zero, returning 0 if the denominator is zero or negative.

\subsection{Data Loading}

The script loads data from \texttt{batch\_summary.csv} using the scenario ID, automatically converting numeric strings to floats for calculations.

\end{document}

